% Plan de tesis de licienciatura Diciembre 2005.
% Pablo Listingart

\documentclass[a4paper, 12pt]{article}
\usepackage[spanish]{babel}
%\usepackage{makeidx}
\usepackage[pdftex]{graphicx}

\addtolength{\topmargin}{-2cm}
\addtolength{\textheight}{4cm}
\makeindex

\begin{document}
\DeclareGraphicsExtensions{.jpg,.pdf,.mps,.png}

\title {
\begin{figure}[h]
\begin{center}
\includegraphics[keepaspectratio, width=3cm]{escudo}
\end{center}
\end{figure}
Universidad de Buenos Aires\\
Facultad de Ciencias Exactas y Naturales\\
Departamento de Computaci\'on\\
\vspace{10mm}
\Huge\textbf{Computaci\'on de alto rendimiento en el estudio de reacciones qu\'imicas en soluci\'on
a trav\'es de m\'etodos cl\'asicos y cu\'anticos.}
\vspace{10mm}
}

\author{por \\
    Pablo Listingart\\
        \\
        Director de Tesis\\
        Lic. Esteban Mocskos\\
        \\
        Propuesta de tesis para optar al grado de \\
        Licenciado en Ciencias de la Computaci\'on\\
        \\
        }
\date{Diciembre de 2005}

\maketitle

\pagebreak

\section*{Propuesta}

Construcci\'on de un software paralelo de simulaci\'on num\'erica mediante el an\'alisis y posterior modificaci\'on de un programa serial preexistente y para su implementaci\'on en un cluster Beowulf con el objeto de estudiar y predecir el comportamiento de reacciones qu\'imicas en soluci\'on.

La combinaci\'on del hardware y software mencionado, y el an\'alisis de distintas estrategias de paralelizaci\'on (incluyendo diversas modalidades de pasaje de mensajes) permitir\'an la resoluci\'on de problemas m\'as realistas y previamente inabordables, obteniendo un algoritmo escalable.


\section*{Objetivos}

\begin{itemize}
\item Desarrollar un programa paralelo modificando el algoritmo de c\'alculo actual utilizando el paradigma de pasaje de mensajes ($Message Passing$) que resuelva num\'ericamente un esquema h\'ibrido cu\'antico-cl\'asico (QM-MM) en el cual el solvente es tratado cl\'asicamente \cite{Nemukhin2002, Kesavan2003, lebrero2002}.

\item Comparar los resultados num\'ericos obtenidos frente a los obtenidos con el programa serial bajo las mismas condiciones y par\'ametros.

\item Incorporar herramientas para realizar el an\'alisis de los datos obtenidos.

\item Utilizar el programa modificado a tama\~nos de problemas m\'as realistas y previamente inabordables.

\end{itemize}

\section*{Elementos a utilizar}

\begin{itemize}
\item Sistema Operativo: Linux, Windows.

\item Lenguaje de programaci\'on: Fortran 77.

\item M\'aquina con suficiente memoria y poder de c\'alculo para poder realizar pruebas iniciales.

\item C\'odigo fuente de la versi\'on serial que soluciona el modelo.

\item Cluster tipo Beowulf y diversas librer\'ias de $Message Passing$ (\texttt{MPICH}\cite{MPICH}, \texttt{LAM}\cite{LAM_MPI}).

\end{itemize}

\section*{Plan inicial}

\begin{enumerate}
\item Interiorizarse de la realidad qu\'imica que ocurre durante un experimento en soluci\'on a trav\'es de m\'etodos cl\'asicos y cu\'anticos.

\item Analizar el programa serial utilizado para la resoluci\'on de sistemas qu\'imicos con el fin de entender los detalles de implementaci\'on del programa existente.

\item Determinar las t\'ecnicas y estrategias posibles de paralelizaci\'on.

\item Realizar la paralelizaci\'on del programa serial, modificando el algoritmo actual a fin de lograr una mejor performance.

\item Realizar experimentos para el estudio de reacciones qu\'imicas en soluci\'on a trav\'es de m\'etodos cl\'asicos y cu\'anticos.

\item Analizar los resultados y realizar pruebas de performance.

\end{enumerate}

\section*{Bibliograf\'ia}

\begin{thebibliography}{9}

\bibitem{Nemukhin2002} A. V. Nemukhin, I. A. Topol, B. L. Grigorenko y S. K. Burt, On the Origin of Potential Barrier for the Reaction OH$^-$ + CO$_2$ $\rightarrow$ HCO$_3$ in Water: Studies by Using Continuum and Cluster Solvation
Methods(2002)

\bibitem{Kesavan2003} L. S. Devi-Kesavan y J. Gao, Combined QM/MM Study of the Mechanism and Kinetic Isotope Effect
of the Nucleophilic Substitution Reaction in Haloalkane
Dehalogenase (2003)

\bibitem{lebrero2002} M. C. Gonz\'alez Lebrero, D. E. Bikiel, M. D. Elola, D. A. Estrin y A. E. Roitberg,
Solvent-induced symmetry breaking of nitrate ion in aqueous clusters: A quantum-classical simulation study.

\bibitem{mpich-install} William D. Gropp and Ewing Lusk, Installation Guide for {\tt mpich}, a Portable Implementation of
{MPI}(1996)

\bibitem{LAM_MPI} Local Area Multicomputer - MPI Parallel Computing Environment, http://www.osc.edu/lam.html

\bibitem{MPICH} MPICH(A Portable MPI Implementation), http://www.mcs.anl.gov/mpi/mpich

\bibitem{Pacheco96} Peter S. Pacheco, Parallel programming with MPI(1996)

\bibitem{Gropp94} William Gropp, Ewing Lusk y Anthony Skjellum, Using MPI: Portable Parallel Programming with the Message-Passing
Interface(1994)

\bibitem{LSC} Laboratorio de Sistemas Complejos, Departamento de Computaci\'on, Facultad de Ciencias Exactas y Naturales, Universidad de Buenos Aires, http://lsc.dc.uba.ar

\bibitem{Amdahl} G.M. Amdahl, Validity of the single-processor approach to achieving large scale computing capabilities(1967)

\bibitem {Gustafson88}John L. Gustafson, Reevaluating Amdahl's Law.(1988)

\end{thebibliography}


\end{document}
